\section{Weak nonlinearity and multiple-scale bookkeeping}
\label{sec:asymptotics}
\subsection{Two small quantities, with a chosen relationship}
Weak nonlinearity means $\varepsilon=k_0a\ll1$. Slow modulation means
$\delta=\Delta k/k_0\ll1$. They are distinct requirements. Dysthe balances
them by taking $\delta=O(\varepsilon)$. A small amplitude alone does not
make an arbitrary spectrum narrowband.

After choosing carrier units, display the scales as
\[
 X=\varepsilon(x-c_gt),\qquad Y=\varepsilon y,\qquad T=\varepsilon^2t.
\]
If $A=\varepsilon\alpha(X,Y,T)$, the chain rule gives
\begin{align*}
 A_x&=\varepsilon^2\alpha_X,& A_{xx}&=\varepsilon^3\alpha_{XX},\\
 A_t&=-c_g\varepsilon^2\alpha_X+\varepsilon^3\alpha_T,&
 (\partial_t+c_g\partial_x)A&=\varepsilon^3\alpha_T.
\end{align*}
The transport cancellation is essential: $A_t$ alone can be second order,
while the derivative following the group is third order.

\begin{definition}[Formal order]
Here $f=O(\varepsilon^m)$ means that the scaled coefficient of
$\varepsilon^m$ is treated as bounded in the chosen slow variables.
Products add orders. A derivative acting on a slow profile adds a factor
of $\varepsilon$, apart from the slower group-frame time scale.
\end{definition}
A rigorous bound would specify the domain, norm, derivative bounds, and
time interval, and prove uniform estimates. The paper's order statements
are formal bookkeeping, not an error theorem for arbitrary initial data.

\subsection{Distinguish slow derivatives from fast derivatives}
For the full oscillatory field,
\[
 \partial_x(Ae^{i\theta})=(A_x+ik_0A)e^{i\theta}.
\]
The carrier derivative $ik_0A$ does not raise the slow order. The statement
after (2.19) that a derivative adds an order concerns derivatives of the
slow coefficients. Carrier derivatives have already been extracted.

The vertical mean-flow scale is slow as well. Its depth dependence is
$e^{Kz}$ with $K=O(\varepsilon)$ in carrier units, whereas the carrier
has $e^z$ dependence.

\begin{center}
\begin{tabular}{@{}ll@{}}
\toprule Quantity in carrier units & Formal order\\\midrule
$A$ & $\varepsilon$\\
$A_x,A_y$ & $\varepsilon^2$\\
$\Dg A$, $\Dg=\partial_t+\tfrac12\partial_x$ & $\varepsilon^3$\\
$A_{xx},A_{yy},A|A|^2$ & $\varepsilon^3$\\
$A_{xxx},A_{xyy},|A|^2A_x,A^2A_x^*$ & $\varepsilon^4$\\
$\Ph$ & $\varepsilon^2$\\
$\Ph_x,\Ph_z,\Zh$ & $\varepsilon^3$\\
$A\Ph_x,A\Ph_z$ & $\varepsilon^4$\\\bottomrule
\end{tabular}
\end{center}

NLS keeps the third-order balance; Dysthe also keeps fourth-order terms.
``Fourth order'' refers to the perturbation expansion, not a fourth
spatial derivative or a quartic nonlinearity.

\subsection{Coefficient matching and resonance}
Insert the harmonic expansion into the boundary conditions. Collect the
coefficient of each $e^{im\theta}$ and the slow mean $m=0$ separately.
Higher harmonics are often determined by forced linear problems in terms
of $A$. They are \emph{bound}: their phase follows the carrier rather than
an independent linear frequency.

The fundamental $m=1$ is special. The leading linear operator already has
the carrier as a homogeneous solution; forcing at that phase can resonate.
The requirement that the next correction remain bounded supplies the
envelope evolution equation.

For comparison, $u_{tt}+u=\varepsilon\cos t$ has a particular term
$\tfrac12\varepsilon t\sin t$, which grows at long times. Allowing the
leading amplitude and phase to evolve slowly removes this secular forcing
from the next correction. This principle produces NLS and its next correction.

The cubic product in Section~\ref{sec:waves} explains the structure
$|A|^2A$ but not its water-wave coefficient. That coefficient comes from
the full boundary conditions, lower harmonics, and normalization.

\subsection{Order reduction}
An intermediate equation contains $A_{tt}$. Leading transport gives
$A_t\approx-\tfrac12A_x$, hence $A_{tt}\approx\tfrac14A_{xx}$.
For fourth-order accuracy one must also differentiate the third-order NLS
correction. Appendix~\ref{sec:reduction} shows this calculation.
Such substitutions are justified only when their errors lie beyond the
target order; they are not identities for arbitrary functions.

\paperlink{Read the order statements around (2.6), (2.12)--(2.17), and
after (2.19). Identify which derivatives act on slow coefficients.}
\checkpoint{For $A=\varepsilon\alpha(\varepsilon x)$, check that
$|A|^2A_x$ is fourth order and $|A|^2A$ third order.}
\clearpage
